% Nagel critique subsection (issue #17), \input by report_kmz.tex.
% Generated artifacts consumed here:
%   \input{\PathToOutput/nagel_comparison_table.tex}   (from src/table_nagel.py)
%   \PathToOutput/figure_nagel.png
% No hand-typed statistics: prose numbers come from report_values.tex macros.

\subsection{Is the virtue of complexity just momentum? Nagel's critique}

\citet{nagel2025seemingly} argues that when the number of random features $P$
greatly exceeds the training window $T$, the ridgeless RFF forecast collapses
to a Gaussian-kernel regression: a similarity-weighted average of the $T=12$
training-window returns. Because the standardized predictors are persistent,
similarity is approximately recency, and kernel distances scale with predictor
volatility, so the strategy is, mechanically, recency-weighted and
volatility-timed momentum. We run his three tests on our own pipeline at the
anchor configuration (the ridgeless model at complexity $c=1{,}000$) and
report what we find, whichever side it favors.

Table~\ref{tab:nagel} scores our replicated VoC strategy against Nagel's
hand-built momentum benchmark, which places linearly declining weights on the
trailing twelve standardized returns, on identical out-of-sample months and
identical metrics (the same performance module used everywhere else in the
project). The benchmark's Sharpe ratio of \NagelBenchSharpe{} sits next to the
VoC row's \NagelVoCSharpe: a strategy with no machine learning in it earns
almost the same risk-adjusted return. One convention note: the VoC row scores
the across-seed mean forecast path, which denoises the forecast slightly, so
it sits a touch above the average-the-statistics value (\MktSharpe) that the
replication reports elsewhere; the comparison within the table is
apples-to-apples.

\begin{table}[htbp]\centering
\input{\PathToOutput/nagel_comparison_table.tex}
\caption{The virtue of complexity is largely spanned by simple momentum. The
VoC strategy and Nagel's recency-weighted momentum benchmark, scored on the
same out-of-sample months with the same metric definitions, deliver similar
performance, and the spanning regression of the VoC strategy on the benchmark
and the static market leaves an alpha with a $t$-statistic of
\NagelSpanAlphaT, below conventional significance.}
\label{tab:nagel}
\end{table}

Figure~\ref{fig:nagel} anatomizes the VoC forecast by regressing it on the
trailing twelve standardized returns; the estimated lag weights are shown
against the benchmark's declining profile. The regression explains only
$R^2 = \NagelAnatomyRTwo$ of the forecast's variation, so the forecast
contains more than recency-weighted momentum, even though its tradable
performance is largely spanned by it.

\begin{figure}[htbp]\centering
\includegraphics[width=0.8\textwidth]{\PathToOutput/figure_nagel.png}
\caption{The VoC forecast is momentum-tilted but not momentum. Estimated lag
weights of the VoC forecast on the trailing twelve standardized returns
(regression $R^2 = \NagelAnatomyRTwo$) against the benchmark's linearly
declining weights: the profiles share the recency tilt, but most of the
forecast's variation is left unexplained.}
\label{fig:nagel}
\end{figure}

\paragraph{The authors' reply.} \citet[Section~4.1]{kelly2025reply} respond
that the momentum reading is a feature, not an artifact: the volatility
standardization is economically motivated, and a reduction toward a
recency-weighted average of past returns is a property any well-specified
nonparametric return predictor would share, not evidence that the
out-of-sample gains are spurious. We present both sides and let the
replicated numbers speak: the performance is real but a much simpler
strategy captures most of it.
